\noindent
\begin{minipage}[t]{0.26\textwidth}
    % Dóng hàng chính xác theo tọa độ y = 635.7pt - 680.4pt (dóng ngang với đẳng thức của Lời giải Ví dụ 5)
    \vspace*{565pt}
    \raggedright
    {\small\textsf{\color{stewartcyan}Lưu ý rằng $\sin 7x \neq 7\sin x$.}}
\end{minipage}%
\hfill
\begin{minipage}[t]{0.71\textwidth}
    \noindent
    (Trong Phụ lục F, bất đẳng thức $\theta \leqslant \tan \theta$ được chứng minh trực tiếp từ định nghĩa độ dài của một cung mà không cần dựa vào trực giác hình học như chúng ta đã làm ở đây.) Do đó, ta có
    \[
    \theta < \frac{\sin\theta}{\cos\theta}
    \]
    \noindent
    \makebox[0pt][l]{nên}%
    \hfill
    $\displaystyle \cos\theta < \frac{\sin\theta}{\theta} < 1$%
    \hfill\null

    \vspace{6pt}
    \noindent
    Ta biết rằng $\lim_{\theta \to 0} 1 = 1$ và $\lim_{\theta \to 0}\cos\theta = 1$, nên theo Định lý kẹp, ta có
    \[
    \lim_{\theta \to 0^+} \frac{\sin\theta}{\theta} = 1 \qquad {\color{stewartcyan}(0 < \theta < \pi/2)}
    \]
    Nhưng hàm số $(\sin\theta)/\theta$ là một hàm chẵn, do đó các giới hạn bên phải và bên trái của nó phải bằng nhau. Vì vậy, ta có
    \[
    \lim_{\theta \to 0} \frac{\sin\theta}{\theta} = 1
    \]
    \noindent
    như vậy chúng ta đã chứng minh được Công thức 5.\hfill{\color{stewartred}$\blacksquare$}

    \vspace{8pt}
    Giới hạn đặc biệt đầu tiên mà chúng ta xem xét liên quan đến hàm sin. Giới hạn đặc biệt tiếp theo sau đây liên quan đến cosin.

    \vspace{8pt}
    \noindent
    \begin{minipage}{\linewidth}
        \centering
        \makebox[0pt][r]{%
            \tcbox[colframe=stewartred, colback=white, boxrule=1pt, sharp corners, left=4pt, right=4pt, top=2pt, bottom=2pt]{%
                \bfseries\color{stewartred}6%
            }\hspace{3.0cm}%
        }%
        \tcbox[colframe=stewartred, colback=white, boxrule=1pt, sharp corners, left=24pt, right=24pt, top=6pt, bottom=6pt]{%
            $\displaystyle \lim_{\theta \to 0} \frac{\cos\theta - 1}{\theta} = 0$%
        }%
    \end{minipage}

    \vspace{10pt}
    \noindent
    \textbf{\textsf{\color{stewartred}CHỨNG MINH}}\quad Ta nhân cả tử số và mẫu số với $\cos\theta + 1$ để đưa hàm số về dạng mà ta có thể áp dụng các giới hạn đã biết.
    \begin{align*}
    \lim_{\theta \to 0} \frac{\cos\theta - 1}{\theta}
    &= \lim_{\theta \to 0} \left( \frac{\cos\theta - 1}{\theta} \cdot \frac{\cos\theta + 1}{\cos\theta + 1} \right) = \lim_{\theta \to 0} \frac{\cos^2\theta - 1}{\theta(\cos\theta + 1)} \\[4pt]
    &= \lim_{\theta \to 0} \frac{-\sin^2\theta}{\theta(\cos\theta + 1)} = -\lim_{\theta \to 0} \left( \frac{\sin\theta}{\theta} \cdot \frac{\sin\theta}{\cos\theta + 1} \right) \\[4pt]
    &= -\lim_{\theta \to 0} \frac{\sin\theta}{\theta} \cdot \lim_{\theta \to 0} \frac{\sin\theta}{\cos\theta + 1} \\[4pt]
    &= -1 \cdot \left( \frac{0}{1 + 1} \right) = 0 \qquad \text{\color{stewartcyan}(theo Công thức 5)} \tag*{{\color{stewartred}$\blacksquare$}}
    \end{align*}

    \vspace{10pt}
    \noindent
    \textbf{\textsf{\color{stewartcyan}VÍ DỤ 5}}\quad Tìm $\displaystyle\lim_{x \to 0} \frac{\sin 7x}{4x}$.

    \vspace{6pt}
    \noindent
    \textbf{\textsf{\color{stewartcyan}LỜI GIẢI}}\quad Để áp dụng Công thức 5, trước tiên ta viết lại hàm số bằng cách nhân và chia cho 7:
    \[
    \frac{\sin 7x}{4x} = \frac{7}{4} \left( \frac{\sin 7x}{7x} \right)
    \]
\end{minipage}
